%%%%%%%%%%%%%%%%%%%%%%%%%%%%%%%%%%%%%%%%%%%%%%%%%%%%%%%%%%%%%%%%%%%%%%%%%%%%%%%
% Curriculum Vitae -- Birmjune Kim
%
% Compile with pdflatex (twice, so cross-references and page breaks settle):
%     pdflatex cv.tex && pdflatex cv.tex
% The resulting cv.pdf is the file the site links to (_config.yml: cv_pdf).
%
% This file is NOT generated from _data/cv.yml. When you update one, update the
% other by hand -- the sections below are kept in the same order as cv.yml.
%%%%%%%%%%%%%%%%%%%%%%%%%%%%%%%%%%%%%%%%%%%%%%%%%%%%%%%%%%%%%%%%%%%%%%%%%%%%%%%

\documentclass[letterpaper,10pt]{article}

% Loaded before fontenc so the T1 fonts get CMaps: without it the fi/ffl
% ligatures come out as garbage when the PDF text is copied or parsed.
\usepackage{cmap}
\usepackage[T1]{fontenc}
\usepackage[utf8]{inputenc}
\usepackage{lmodern}
% Bold resolves to Latin Modern Roman Demi (a semibold) rather than the heavier
% bold-extended, which is what gives headings their weight in the reference CV.
\renewcommand{\bfdefault}{sb}
\usepackage[letterpaper,top=0.5in,bottom=0.45in,left=0.7in,right=0.7in]{geometry}
\usepackage{enumitem}
\usepackage[hidelinks]{hyperref}

\pagestyle{empty}
\setlength{\parindent}{0pt}
\raggedbottom

%%%%%%%%%%%%%%%%%%%%%%%%%%%%%%%%%%% Layout macros %%%%%%%%%%%%%%%%%%%%%%%%%%%%%%

% Section heading: \Large semibold serif, flush with the left margin, underlined
% by a rule spanning the full text width.
% \smash plus a fixed-height zero-depth rule gives every heading an identical
% line box, so descenders ("Experience", "Projects") cannot push the rule -- and
% the text below it -- further down than in the other sections.
\newcommand{\heading}[1]{%
  \par\vspace{6.5pt}%
  {\Large\bfseries\rule{0pt}{10pt}\smash{#1}}\par
  \vspace{4.5pt}%
  \hrule height 0.4pt
  \vspace{5pt}%
}

% Everything under a heading sits in this indented block.
\newenvironment{resbody}%
  {\begin{list}{}{%
     \setlength{\leftmargin}{0.175in}%
     \setlength{\rightmargin}{0pt}%
     \setlength{\topsep}{2pt}%
     \setlength{\parsep}{0pt}%
     \setlength{\itemsep}{0pt}%
   }\item[]}%
  % Close on a strut of known depth (and take its line back out again) so the
  % gap to the next heading barely depends on whether the section ended in a
  % paragraph or in a bullet list.
  {\par\strut\vspace{-\baselineskip}\end{list}}

% One line, left part and right part pushed to the margins. Ending the line with
% \par instead of \\ (and zeroing \parfillskip so \hfill survives) keeps LaTeX
% from adding a stray empty line before a following bullet list.
\newcommand{\rline}[2]{{\setlength{\parfillskip}{0pt}\strut #1\hfill #2\par}}

% One entry: bold title on the left, date flush right.
\newcommand{\entry}[2]{\rline{\textbf{#1}}{#2}}

% Second line of an entry: role on the left, location or result flush right.
\newcommand{\subentry}[2]{\vspace{2pt}\rline{\textit{#1}}{\textit{#2}}}

% Flat "item ......... year" line, used for honors. The extra leading keeps the
% long single-line list from reading as a solid block.
\newcommand{\award}[2]{\rline{#1}{#2}\vspace{1.6pt}}

% Space between two entries of the same section.
\newcommand{\entrygap}{\vspace{3pt}}

% Dotted bullets that hang into the left indent, as in the reference layout.
\newlist{points}{itemize}{1}
% itemsep matches the extra leading \award adds, so a bullet list and the honors
% list breathe the same amount between entries.
\setlist[points]{topsep=2.5pt, partopsep=0pt, parsep=0pt, itemsep=1.6pt,
                 leftmargin=0pt, labelwidth=6pt, labelsep=2.5pt, label=$\cdot$}

%%%%%%%%%%%%%%%%%%%%%%%%%%%%%%%%%%%%%% Header %%%%%%%%%%%%%%%%%%%%%%%%%%%%%%%%%%

\begin{document}

\begin{center}
  {\huge\bfseries Birmjune Kim}

  \vspace{6pt}
  \href{mailto:johnbie0627@gmail.com}{johnbie0627@gmail.com}$\diamond$%
  \href{https://birmjune.github.io}{Personal Website}$\diamond$%
  \href{https://github.com/Birmjune}{GitHub}
\end{center}

%%%%%%%%%%%%%%%%%%%%%%%%%%%%%%%%%%%% Education %%%%%%%%%%%%%%%%%%%%%%%%%%%%%%%%%

\heading{Education}
\begin{resbody}
  \strut\textbf{Seoul National University (SNU)}, Seoul, Republic of Korea\par
  \begin{points}
    \item B.S. in Computer Science and Engineering \hfill Mar 2024 - Present
    \item Kim Jaechul AI Class \hfill Sep 2026 - Present
    \item Cumulative GPA: 4.24 / 4.3~$\cdot$~Major GPA: 4.27 / 4.3
  \end{points}

  \entrygap
  \rline{\textbf{Seoul Science High School}, Seoul, Republic of Korea}{Mar 2021 - Feb 2024}
\end{resbody}

%%%%%%%%%%%%%%%%%%%%%%%%%%%%%%% Research interests %%%%%%%%%%%%%%%%%%%%%%%%%%%%%

\heading{Research Interests}
\begin{resbody}
  \begin{points}[itemsep=5pt]
    \item \textbf{Efficient AI.} Making efficient AI models, across inference time,
      compression and quantization, AI compilers, and hardware acceleration; currently
      applied to foundation models and time series modeling.
    \item \textbf{AI in the Physical World.} Putting those efficient models to work in the
      physical world (like robotics), where latency, memory, and power are real constraints.
  \end{points}
\end{resbody}

%%%%%%%%%%%%%%%%%%%%%%%%%%%%%%%%%%% Experience %%%%%%%%%%%%%%%%%%%%%%%%%%%%%%%%%

\heading{Experience}
\begin{resbody}
  \entry{Machine Learning Lab, Seoul National University}{Dec 2025 - Present}
  \subentry{Research Intern (Advisor: Prof. Hyun Oh Song)}{Seoul, South Korea}
  \begin{points}
    \item Researching a dynamic programming based dynamic decoding process for time series
      foundation models.
  \end{points}
\end{resbody}

%%%%%%%%%%%%%%%%%%%%%%%%%%%%%%%%%%%% Projects %%%%%%%%%%%%%%%%%%%%%%%%%%%%%%%%%%

\heading{Projects}
\begin{resbody}
  \entry{\href{https://github.com/Birmjune/2026_SNU_AI_Challenge_DeepRed}{SNU AI Challenge 2026: Storyline Frame Ordering}}{Jul 2026 - Aug 2026}
  \begin{points}
    \item Given a caption describing a storyline and four shuffled frames from that story,
      predict each frame's position in the correct temporal order.
    \item Finalist, one of 12 teams out of 206; placed 10th overall.
  \end{points}

  \entrygap
  \entry{\href{https://github.com/Birmjune/Image2Answer-4optionMCQ}{Image2Answer: 4-Option MCQ Solver}}{Aug 2025}
  \begin{points}
    \item A vision-language pipeline that reads a question image and predicts the correct
      choice among the four options.
    \item Samsung Collegiate Programming Challenge (AI track): one of 41 finalists out of
      1,483 entrants.
  \end{points}

  \entrygap
  \entry{\href{https://github.com/Birmjune/Easy-ARC-AGI-solver}{Easy ARC-AGI Challenge Solver}}{Spring 2025}
  \begin{points}
    \item Infers the rule behind a few example grid pairs and applies it to an unseen input,
      for grids up to $10 \times 10$.
    \item Private score 63/100.
  \end{points}
\end{resbody}

%%%%%%%%%%%%%%%%%%%%%%%%%%%%%%% Honors and awards %%%%%%%%%%%%%%%%%%%%%%%%%%%%%%

\heading{Honors and Awards}
\begin{resbody}
  \award{Finalist, 10th of 206 teams, SNU AI Challenge 2026}{2026}
  \award{One of 41 finalists (1,483 entrants), Samsung Collegiate Programming Challenge (AI track)}{2025}
  \award{National Science \& Technology Scholarship (Full Scholarship)}{2024}
  \award{8th of 275 pairs, Simon Marais Mathematics Competition (Pairs division, East)}{2024}
  \award{Best Award, Graduation Thesis in Mathematics, Seoul Science High School}{2023}
  \award{Gold Medal, Korean Mathematical Olympiad (KMO), high school division}{2022}
  \award{Best Award, R\&E Research Competition in Mathematics, Seoul Science High School}{2022}
\end{resbody}

%%%%%%%%%%%%%%%%%%%%%%%%%%%%%%% Skills and languages %%%%%%%%%%%%%%%%%%%%%%%%%%%

\heading{Skills}
\begin{resbody}
  \strut\textbf{Programming.} Python, PyTorch, C / C++, Shell \quad
  \textbf{Languages.} Korean (native), English
\end{resbody}

%%%%%%%%%%%%%%%%%%%%%%%%%%%%%%%%%%%%%% Footer %%%%%%%%%%%%%%%%%%%%%%%%%%%%%%%%%%

\vfill
\begin{center}
  {\small Last updated: August 10, 2026}
\end{center}

\end{document}
